%%%%%%%%%%%%%%%%%%%%%%%%%%%%%%%%%%%%%%%%%%%%%%%%%%%%%%%%%%%%%%%%%%%%%%%%%%%%%%
\section{Background}
%%%%%%%%%%%%%%%%%%%%%%%%%%%%%%%%%%%%%%%%%%%%%%%%%%%%%%%%%%%%%%%%%%%%%%%%%%%%%%

% - . - . - . - . - . - . - . - . - . - . - . - . - . - . - . - . - . - . - .%
\subsection{Segment Graph}
% - . - . - . - . - . - . - . - . - . - . - . - . - . - . - . - . - . - . - .%
- Similarly to \cite{Matar18}, Taskgrind reason on a \emph{segment graph}
- nodes = non-divisible sequence of instructions
- there is a path between Ni and Nj if a synchronization imposes the ordering of the two sequences of instructions (said differently: if a path exists, then 'Ni' happens-before 'Nj').


% - . - . - . - . - . - . - . - . - . - . - . - . - . - . - . - . - . - . - .%
\subsection{Valgrind}
% - . - . - . - . - . - . - . - . - . - . - . - . - . - . - . - . - . - . - .%
- Valgrind tool = core + a tool plugin
- JIT re-compilation of code block from binary code to an IR (VEX)
- the IR tree is passed to the tool that can inject IR instructions for instrumenting
- in addition, valgrind provides many utilities
	- Client requests
	- Function replacement, for instance to wrap the memory allocator
	- ...
	
from [Valgrind]- R3: Instrument read/write instructions. Valgrind supports this
requirement—all reads and writes of registers and memory are vis-
ible in the IR and instrumentable. The IR’s load/store nature makes
instrumentation of memory accesses particularly easy

%%%%%%%%%%%%%%%%%%%%%%%%%%%%%%%%%%%%%%%%%%%%%%%%%%%%%%%%%%%%%%%%%%%%%%%%%%%%%%
\section{Taskgrind: a Valgrind tool for determinacy race detection}
%%%%%%%%%%%%%%%%%%%%%%%%%%%%%%%%%%%%%%%%%%%%%%%%%%%%%%%%%%%%%%%%%%%%%%%%%%%%%%

Scheme showing the Application (client) / Taskgrind (the tool) / Analyzer / Valgrind (the core) / the PM tool (OMPT, Qthreads...)
	- in application : any binary program
	- in Valgrind : utilities R3, R8, ...
	- in Taskgrind : segment graph building, memory accesses instrumentation, allocator replacement, TLS detection... then detect determinacy race, report errors
	- in PM Tool (OMPT) : callbacks forwarding segments to Taskgrind
		- Note: Taskgrind could use function replacement to wrap any parallel libraries and re-implement OMPT internally, for instance
		- but interfaces exists, more convenient and portable to use them



- Taskgrind builds a segment graph during the execution, and keep track of 'the current segment'
- Accesses are accumulated in compact interval trees (one for load, one for stores)
(self-balancing red-black tree)	
- Finding determinacy race given this graph becomes trivial - O(n²) with 'n' the number of nodes



% - . - . - . - . - . - . - . - . - . - . - . - . - . - . - . - . - . - . - .%
\subsection{Recording Parallel Execution}
% - . - . - . - . - . - . - . - . - . - . - . - . - . - . - . - . - . - . - .%

\paragraph{OpenMP}
- OMPT tool
- Convert OMPT callbacks to a Taskgrind segment graph
- create(T) - the executing thread creates the task 'T'
- schedule(T) - the executing thread schedules the task 'T'
- depend(T, list<in|out|outset, addr>) - openmp-like dependences
- the first segment the successor depends on the last segment of the predecessor
- detach
- when omp-fulfill-event(e) is raised, create an empty segment S
- successors of e.T also depends on S
- sync(children|all) - the executing thread must wait for
- all children tasks completion
- all descendant tasks completion

Remark:
- a parallel-for loop chunk could be interpreted as a single-segment ttask - cannot map a chunk creation (work callback) with its schedule (dispatch callback)
- a target
- nowait 	- is a task
- wait 		- is an undefered task
so targetting CPU, it should work

\paragraph{QThreads}

\paragraph{Cilk}

% - . - . - . - . - . - . - . - . - . - . - . - . - . - . - . - . - . - . - .%
\subsection{Collecting Memory Accesses}
% - . - . - . - . - . - . - . - . - . - . - . - . - . - . - . - . - . - . - .%
- Accesses are accumulated in compact interval trees (one for load, one for stores)
(self-balancing red-black tree)	


% - . - . - . - . - . - . - . - . - . - . - . - . - . - . - . - . - . - . - .%
\subsection{Analyzing Instrumented Execution}
% - . - . - . - . - . - . - . - . - . - . - . - . - . - . - . - . - . - . - .%

Taskgrind not only represent the execution 


\paragraph{Determinacy Race}
Finding determinacy race given this graph becomes trivial - O(n²) with 'n' the number of segments

foreach segment s1
	foreach segment s2 != s1
		if no path exists between s1 and s2
			foreach address of s1.W n (s2.read u s2.write)
				report determinacy race on X from segments s1 and s2

No false-negative, hence detecting semantical errors even if the execution was correct.

\paragraph{???}



%%%%%%%%%%%%%%%%%%%%%%%%%%%%%%%%%%%%%%%%%%%%%%%%%%%%%%%%%%%%%%%%%%%%%%%%%%%%%%
\section{Pitfalls with Heavyweight Instrumentation}
%%%%%%%%%%%%%%%%%%%%%%%%%%%%%%%%%%%%%%%%%%%%%%%%%%%%%%%%%%%%%%%%%%%%%%%%%%%%%%

But, way too many false-positive. We show a few examples
       
       
% - . - . - . - . - . - . - . - . - . - . - . - . - . - . - . - . - . - . - .%
\subsection{Memory Recycling}
% - . - . - . - . - . - . - . - . - . - . - . - . - . - . - . - . - . - . - .%
- can easily wrap memory allocation ; we override the allocator and transform 'free' to a 'no-op'
- problem with LLVM - it has its own allocator for internal data structures that may be recycled (e.g. tasks descriptor)
	- raises many false-positive, Taskgrind currently suppress any determinacy races on memory allocated by this allocator.

% - . - . - . - . - . - . - . - . - . - . - . - . - . - . - . - . - . - . - .%
\subsection{Thread-Local Accesses}
% - . - . - . - . - . - . - . - . - . - . - . - . - . - . - . - . - . - . - .%
- Stack accesses
	- we register the frame pointer
- Thread Local Storage (TLS)~\cite{elf-tls}
	- standard with `-Thread-local` : OK, Valgrind already has some routines to help parsing TCB/DTV
	- non-standard
		- `array[omp-get-thread-num()]` or same with pthreads... can hardly be detected

% - . - . - . - . - . - . - . - . - . - . - . - . - . - . - . - . - . - . - .%
\subsection{Parallel Runtimes Non-Determinacy}
% - . - . - . - . - . - . - . - . - . - . - . - . - . - . - . - . - . - . - .%
Non-determinacy is not always an issue.
For instance, most OpenMP implementation task-scheduling is non-deterministic it-self : given a task-based program and its input, tasks may be scheduled in different partial order and on different cores.
We simply blacklist all the parallel runtime system symbols from instrumentation to avoid reporting such issues.





